\chapter{検証障害と選択バイアス}
\label{ch:method}
%============================================================

\section{問題設定}

第\ref{part:theory}部の枠組みを実企業に適用しようとすると、
比較の出発点に系統的な歪みが入り込む。この歪みは単一の「バイアス」ではなく、
性質の異なる少なくとも六種の問題の混合である。以下ではそれを分解し、
対処可能なものと原理的に不可能なものを区別する。

\section{生存条件づけ：枠組みの循環}
\label{sec:survivorship}

観測可能な企業とは、制約\eqref{eq:constraint}を破らなかった経路の実現値である。
倒産時刻を $\tau=\inf\{t:M(t)<0\}$ とすると、我々が観測するのは
\begin{equation}
\Prob(X \mid \tau > t)
\end{equation}
であって $\Prob(X)$ ではない。そして $\tau$ を決めているのは、まさに推定対象である $\Phi$ である。

\begin{figure}[htbp]
\centering
\begin{tikzpicture}[>=Stealth,node distance=1.2cm]
  \node[draw,rounded corners,minimum width=3.4cm,minimum height=1.0cm,align=center,font=\small]
    (phi) at (0,1.6) {$\Phi$：ビジネスモデル\\[-2pt]\scriptsize 調べたい構造};
  \node[draw,rounded corners,minimum width=3.4cm,minimum height=1.0cm,align=center,font=\small]
    (eps) at (0,-1.6) {$\varepsilon$：衝撃・運\\[-2pt]\scriptsize 外生的な変動};
  \node[draw,very thick,rounded corners,minimum width=3.8cm,minimum height=1.0cm,align=center,font=\small,fill=gray!15]
    (s) at (6.2,0) {$S$：生存・開示\\[-2pt]\scriptsize $S=1$ で条件づけ};
  \draw[->,thick] (phi) -- (s);
  \draw[->,thick] (eps) -- (s);
  \draw[<->,dashed] (phi) to[bend right=45] node[right,font=\scriptsize,align=left,xshift=2pt] {条件づけで生じる\\見かけの相関} (eps);
\end{tikzpicture}
\caption{生存 $S$ は $\Phi$ と $\varepsilon$ の共通の帰結（合流点）である。$S$ で条件づけると両者に見かけの相関が生じる。}
\label{fig:dag}
\end{figure}

$S$ は $\Phi$ と外生的衝撃 $\varepsilon$ の共通の帰結、すなわち合流点である。
合流点で条件づけると、本来独立な二つの原因の間に相関が生じる（Berkson のパラドックス）。

\subsection{歪みの向きは条件づけの対象で反転する}

これは実務上重要な区別である。

\paragraph{生存で条件づける場合。}
生存企業のうち、不利な $\Phi$ を持ちながら存続した企業は $\varepsilon$ が良好だった企業である。
したがって生存サンプルでは $\Phi$ の劣位が $\varepsilon$ の優位で相殺されて見え、
\textbf{構造の効果は過小推定される}。$\CCC>0$ であることの不利は、
存続企業のみを見ていると実際より小さく見える。

\paragraph{成功（有名さ）で条件づける場合。}
こちらは事情が異なる。成功企業のみを収集すると、
同じ $\Phi$ を持ちながら失敗した企業が標本から完全に消えるため、対照群が存在しない。
過小推定ではなく\textbf{そもそも推定が不可能}であり、
$\Phi$ と $\varepsilon$ を分離する情報が標本内に存在しない。

両者を同一の「生存バイアス」として扱うと、この差が見えなくなる。

\section{開示選択：調べたい軸との相関}
\label{sec:disclosure}

上場・開示は無作為な選択ではない。式\eqref{eq:gstar}より、
$\CCC<0$ の企業は $g^\star$ が存在せず成長を自己金融できるため、
外部資本市場に出る必要が原理的に薄い。
逆に $\CCC$ が大きく正の企業は外部資本を要し、上場に押し出されやすい。

すなわち\textbf{選択確率が、まさに推定対象の関数になっている}。
選択確率 $\Prob(S_{\mathrm{disc}}=1 \mid \CCC)$ は $\CCC$ の増加関数となる傾向を持つ。
したがって上場企業データベースで観測される $\CCC$ の分布は、母集団の分布から系統的に、
かつ予測可能な方向にずれている。

\section{年齢：バイアスではなく定義域の問題}
\label{sec:age}

設立初日の企業と数年経過した企業を比較できないという問題は、性質が異なる。
式\eqref{eq:ccc}は定常状態を仮定して導かれており、
$r\approx 0$ ではゼロ除算となって定義されない。

したがってこれはバイアスではなく\textbf{枠組みの適用範囲の問題}である。
対処法も異なる。バイアスなら補正するが、定義域の問題であれば宣言するしかない。
具体的には次のいずれかを取る。

\begin{enumerate}[label=(\arabic*)]
\item $r$ が安定した領域に対象を限定し、閾値を明示する。
\item 過渡状態を別の変数系で明示的にモデル化する。
\item 分析単位を企業から $\Phi$ に下げ、$\Phi$ ごとの成熟度で判定する（第\ref{sec:unit}節）。
\end{enumerate}

補正係数を導入して定常状態の式を無理に適用することは、最も避けるべき対処である。

\section{年齢・時点・コホートの識別不能性}
\label{sec:apc}

企業を観測したとき、三つの時間が同時に効いている。設立年をコホート $c$、
観測年を $p$、年齢を $a$ とすると、
\begin{equation}
a = p - c
\label{eq:apc-main}
\end{equation}
が恒等的に成立する。式\eqref{eq:apc-main}により三変数は線形従属であり、
\textbf{三つの効果を同時に識別することは、外生的な制約を追加しない限り不可能}である。
これは Age--Period--Cohort 問題として人口学および疫学で長く議論されてきたが、
一般的な解決は存在しない。

したがって「ある年代のある業種の $\CCC$ が良好なのは、$\Phi$ によるのか、若さによるのか、
当時の金融環境によるのか」という問いには、データのみでは答えられない。
いずれかの効果をゼロと置くなどの識別仮定を明示し、
結論がその仮定にどれだけ依存するかを示すほかない。

\section{事後的物語化}
\label{sec:narrative}

成功した企業の $\Phi$ の記述それ自体が汚染されている。
構造は成功を説明するように語り直されるため、
$\Phi$ という説明変数の\emph{測定}にバイアスが入る。
これは従属変数側の補正では除去できない。

第\ref{part:theory}部で述べたとおり、$\Phi$ の多くは設計の産物ではなく業態の性質として生じ、
事後に金融的機能が認識されたものである。
一次資料（当時の開示、契約書、報道）と回顧的記述を区別することが最低限の対処になる。

\section{観測フレームの実際}
\label{sec:frame}

日本を対象とする場合の規模感を示す。

\begin{center}
\begin{tabular}{@{}lrl@{}}
\toprule
母集団 & 規模 & 出典・注 \\
\midrule
企業数（中小企業庁ベース） & 約 360 万 & 経済センサスの再編加工。一次産業等を除く \\
企業数（国税庁ベース） & 約 640 万 & 個人事業主約 379 万を含む \\
上場企業 & 3{,}975 社 & 2026年8月27日時点 \\
日常的に言及される企業 & 数百社 & 全体の $10^{-4}$ の桁 \\
\bottomrule
\end{tabular}
\end{center}

上場企業は全体のおよそ 0.1\%、有名企業は 0.01\% 以下である。
「様々な企業」と述べるとき、どの母集団を指しているかが決定的に効く。

\section{対処方針}

対処可能なものと不可能なものを分けて列挙する。

\subsection{フレームを事前に宣言し、前向きに追跡する}

「有名企業を集める」のではなく、「基準時点 $t_0$ に存在した特定業種の全事業者」を定義し、
そこから前方に追跡する。こうすると生存は隠れたフィルタではなく観測される従属変数になり、
第\ref{sec:survivorship}節の合流点条件づけが回避される。

利用可能なフレームとして、日本では経済センサス、
帝国データバンクおよび東京商工リサーチのデータベースがある。
欧州では Orbis が非上場企業を広く含む。
米国は開示義務が緩いため非上場データが相対的に弱く、
この点では日欧のほうが条件が良い。

\subsection{失敗側を明示的に収集する}

倒産は官報および信用調査会社の集計によって観測可能である。
分母を構成できるならば、生存バイアスは推定可能な選択モデルに変わる。

\subsection{分析単位を企業から \texorpdfstring{$\Phi$}{Φ} に下げる}
\label{sec:unit}

これが最も効果の大きい対処である。
第\ref{sec:phi}章の定義より、$\Phi$ は契約レベルの対象であり、企業はその束である。
大規模企業は通常、$\kappa$ の符号すら異なる複数の $\Phi$ を並行して運用している。
企業単位で $\CCC$ を計算することは、これらを混合した平均を見ることに等しい。

分析単位を $\Phi$ に下げることで、
\begin{itemize}
\item 年齢問題が部分的に解消される（成熟企業内の新しい $\Phi$ を扱える）
\item 観測数が増える
\item 「企業」という制度的人工物の影響を減らせる
\end{itemize}

\subsection{産業分類ではなく制約の形でグルーピングする}

産業分類は特定時代の産業観を固定したものであり、国によっても定義が異なる。
第\ref{part:theory}部の枠組みが名指しした構造的特徴、すなわち
\begin{itemize}
\item 容量制約の有無（第\ref{sec:breakage}節）
\item 取引頻度 $\nu$（式\eqref{eq:decay}）
\item 努力の観測可能性（第\ref{sec:info}章）
\item 資本集約度（式\eqref{eq:cash}の $A$ の比重）
\end{itemize}
でグルーピングすれば、業種と時代をまたいだ比較が意味を持つ。
理論を構成した見返りはここにある。

\subsection{除去できないものは符号を明示する}

第\ref{sec:apc}節の識別不能性と第\ref{sec:narrative}節の物語化は除去できない。
ただし方向を明示することはできる。
「この推定値は構造の効果を過小に見積もっている」
「この $\Phi$ の記述は事後的合理化を含む可能性が高い」と述べれば、
結論の頑健性は判断可能になる。

\subsection{有名企業の使用規則}

以上より、次の規則を置く。
\begin{quote}
\textbf{有名企業は例示に用い、推論には用いない。}
\end{quote}
構造の説明として個別企業を挙げることは許容される。
成功例の存在から構造の優位を結論することは許容されない。

%============================================================
